\section{Distribution Power Flow}

\subsection{Problem Form and Runtime Scope}

This chapter covers the distribution backward-forward-sweep solver for radial AC
feeders.  The implemented mathematical model is a current-injection radial load
flow with outer reactive-limit enforcement, distinct from the transmission-side
hybrid Newton formulation documented in Section~\ref{sec:powerflow}.

\subsection{Code Map}

The current implementation is concentrated in:
\begin{itemize}
  \item \texttt{src/analysis/distribution\_power\_flow.cpp} for BFS execution and reporting,
  \item \texttt{src/analysis/topology\_analysis.cpp} for radiality and connectivity checks,
  \item shared model structs in the main information-model layer for buses, branches, DERs, and loads.
\end{itemize}

\subsection{Pseudocode Map}

The relevant code-aligned workflow is Algorithm A14 in
Section~\ref{sec:algorithms}.  The feeder-topology assumptions discussed here
also align with the ONR workflow in Algorithm A13.

\subsection{Current Status}

The BFS solver is implemented and usable for radial feeders, but its model scope
is intentionally narrower than the canonical projection used by the main PF
stack.  It is therefore best read as a specialized feeder-analysis solver rather
than a universal hybrid AC/DC load-flow engine.

\subsection{Scope of the Current DPF Module}

The distribution solver implements a backward-forward sweep (BFS) for radial AC
networks.  Unlike the main PF stack, it does \emph{not} automatically invoke the
canonical projection layer.  Its practical input coverage is therefore:
\begin{itemize}
  \item direct \texttt{ACBus}, \texttt{ACBranch}, \texttt{Generator}, \texttt{Load}, \texttt{Shunt}, \texttt{StaticGenerator}, \texttt{RenewableGen}, \texttt{PVSystem}, \texttt{Storage}, and \texttt{ChargingStation},
  \item hybrid overload only for fixed \texttt{VirtualPowerPlant}, \texttt{Microgrid}, and \texttt{MobileStorage} injections,
  \item no automatic use of \texttt{Transformer2W}, \texttt{Transformer3W}, \texttt{Switch}, \texttt{FlexibleLoad}, \texttt{AsymmetricLoad}, \texttt{Charger}, or optional three-phase models.
\end{itemize}

\subsection{Radial Network Requirement}

The solver assumes a single slack bus and a radial branch topology.  The helper
functions \texttt{is\_connected()} and \texttt{is\_radial()} in
\texttt{topology\_analysis.hpp} are intended to be used before calling the BFS
solver on a candidate feeder model.

\subsection{Net-Injection Model}

For each non-root bus, the BFS routine constructs a complex net demand
\begin{align}
S_i^{net} = P_i^{net} + jQ_i^{net},
\end{align}
with the sign convention ``load minus generation'' rather than ``bus injection
positive.''  In current code form,
\begin{align}
S_i^{net} &= S_i^{load} + S_i^{cs} + S_i^{shunt} - S_i^{gen} - S_i^{sg} - S_i^{ren} - S_i^{pv} - S_i^{stor},
\end{align}
plus optional hybrid injections from VPP, microgrid export, and stationary
mobile storage when the \texttt{HybridPowerSystem} overload is used.

\paragraph{Important implementation detail.}
Even though \texttt{Load} stores ZIP coefficients, the current BFS routine uses
\texttt{Load::p\_mw} and \texttt{Load::q\_mvar} directly as a constant-power
baseline and does not iterate the full ZIP model against bus voltage.  The same
simplification applies to shunts: when included, they are added as fixed complex
power contributions rather than true \(V^2\)-dependent admittance terms.

\subsection{Backward-Forward Sweep Equations}

For branch current and bus voltage updates, the current implementation follows
standard radial BFS logic.

\paragraph{Backward sweep.}
At bus \(j\), the injected current is approximated as
\begin{align}
I_j^{inj} = \overline{\frac{S_j^{net}}{V_j}},
\end{align}
and branch currents are accumulated from children toward the root.

\paragraph{Forward sweep.}
If branch \(b\) connects parent bus \(i\) to child bus \(j\), then
\begin{align}
V_j = V_i - Z_b I_b,
\qquad Z_b = r_b + jx_b.
\end{align}

\paragraph{Convergence metric.}
The stopping criterion is the maximum voltage change:
\begin{align}
\varepsilon = \max_i |V_i^{(k)} - V_i^{(k-1)}|.
\end{align}

\subsection{Outer Loop for Reactive Limits}

When \texttt{DPFOptions::enforce\_q\_limits} is enabled, the solver runs an outer
loop that:
\begin{enumerate}
  \item solves the BFS problem,
  \item estimates generator reactive output at PV buses,
  \item clamps generators that exceed \(Q\) limits,
  \item converts affected buses from PV behavior to PQ behavior,
  \item repeats until no new violations are found or the outer limit is hit.
\end{enumerate}

\subsection{Reported Branch Power and Losses}

After convergence, the module reports sending-end branch power
\begin{align}
S_b^{from} = V_{from(b)}\overline{I_b},
\end{align}
and network loss totals
\begin{align}
P_{loss}^{tot} &= \sum_b |I_b|^2 r_b S_{base}, \\
Q_{loss}^{tot} &= \sum_b |I_b|^2 x_b S_{base}.
\end{align}

\subsection{Relationship to the Main Newton PF}

The BFS implementation is best viewed as a feeder-oriented fast solver for radial
networks.  It is \emph{not} a drop-in replacement for the main projected hybrid
PF path because:
\begin{itemize}
  \item it does not consume the same projected component set,
  \item it does not solve the DC subsystem,
  \item it approximates load and shunt voltage dependence more aggressively,
  \item it uses a radial tree traversal rather than the sparse Jacobian machinery.
\end{itemize}

% ===================================================================
% Three-Phase (abc-domain) Extension
% ===================================================================
\subsection{Three-Phase Power Flow (abc Domain)}
\label{sec:dpf-three-phase}

Two three-phase solvers are provided, both operating in the full abc phase
domain:
\begin{enumerate}
  \item \textbf{Newton--Raphson (NR)} --- the primary solver, suitable for
        general (meshed, multi-source, and hybrid AC/DC) network topologies.
        It builds a $3n\times3n$ admittance matrix and solves polar-form
        mismatch equations via Newton's method.
  \item \textbf{Backward--Forward Sweep (BFS)} --- a legacy solver retained for
        radial single-source distribution feeders.  See
        Section~\ref{sec:dpf-three-phase-bfs}.
\end{enumerate}

\subsubsection{Phase Domain and API}

An enumeration \texttt{PhaseDomain} selects between positive-sequence (default
single-phase) and abc (per-phase) analysis.

The NR solver is invoked through \texttt{solve\_three\_phase\_nr()}, and the
convenience wrapper \texttt{solve\_three\_phase()} (in
\texttt{powerflow/three\_phase.hpp}) delegates to it.
Both accept a \texttt{ThreePhaseACSystem} or the \texttt{HybridPowerSystem}
wrapper (which checks for a populated \texttt{three\_phase\_ac} subsystem).

The legacy BFS solver remains available via
\texttt{solve\_three\_phase\_distribution\_pf()}.

\subsubsection{\texorpdfstring{$3\times3$}{3x3} Series Impedance from Sequence Parameters}

Each \texttt{ThreePhaseACLine} carries positive-sequence $Z_1 = r_1 + jx_1$ and
zero-sequence $Z_0 = r_0 + jx_0$ impedance.  The $3\times3$ abc-domain series
impedance matrix is built via the Fortescue-to-abc transformation:

\begin{align}
  \mathbf{Z}_{abc} &= \mathbf{A}\,
    \operatorname{diag}(Z_0,\, Z_1,\, Z_1)\,
    \mathbf{A}^{-1},
\end{align}
where $\mathbf{A}$ is the symmetrical component transformation matrix
($a = e^{j2\pi/3}$):
\begin{align}
  \mathbf{A} = \frac{1}{3}
  \begin{bmatrix}
    1 & 1 & 1 \\
    1 & a^2 & a \\
    1 & a & a^2
  \end{bmatrix}.
\end{align}
The resulting impedance is a circulant matrix:
\begin{align}
  \mathbf{Z}_{abc} =
  \begin{bmatrix}
    Z_s & Z_m & Z_m \\
    Z_m & Z_s & Z_m \\
    Z_m & Z_m & Z_s
  \end{bmatrix},
  \qquad
  Z_s = \frac{Z_0 + 2Z_1}{3}, \quad
  Z_m = \frac{Z_0 - Z_1}{3}.
\end{align}
When the line is balanced ($Z_0 = Z_1$), the mutual impedance $Z_m = 0$ and the
three phases decouple into independent single-phase circuits, each equivalent to
the positive-sequence model.

\subsubsection{Per-Phase Net Injection}

At each non-slack bus $i$, the three-phase net demand vector is
\begin{align}
  \mathbf{S}_i^{net} =
  \begin{bmatrix} S_i^a \\ S_i^b \\ S_i^c \end{bmatrix}
  =
  \begin{bmatrix}
    P_i^{d,a} + jQ_i^{d,a} \\
    P_i^{d,b} + jQ_i^{d,b} \\
    P_i^{d,c} + jQ_i^{d,c}
  \end{bmatrix}
  - \mathbf{S}_i^{gen},
\end{align}
where bus-level demand $P_i^{d,\phi}$, $Q_i^{d,\phi}$ and three-phase load
table entries are summed per phase.  Generators with per-phase output are
subtracted, using equal $P/3$, $Q/3$ splits across phases when only aggregate
ratings are available.

% -------------------------------------------------------------------
% Three-Phase Newton-Raphson Solver
% -------------------------------------------------------------------
\subsubsection{Three-Phase Newton--Raphson Solver}
\label{sec:dpf-three-phase-nr}

The three-phase Newton--Raphson solver generalises the standard single-phase
polar-form NR power flow to a $3n\times3n$ phase-node system, where $n$ is the
number of buses.

\paragraph{Phase-node indexing.}
Each bus $i\in\{0,\ldots,n{-}1\}$ contributes three phase nodes:
\begin{align}
  r = 3i + \phi, \qquad \phi \in \{0\,(a),\;1\,(b),\;2\,(c)\}.
\end{align}

\paragraph{\texorpdfstring{$3n\times3n$ admittance matrix $\mathbf{Y}^{abc}$}{3n x 3n admittance matrix Yabc}.}
For each three-phase line connecting bus $i$ to bus $j$, the $3\times3$ abc
admittance block is obtained by inverting the circulant impedance:
\begin{align}
  Y_s &= \frac{1/Z_0 + 2/Z_1}{3}, \qquad
  Y_m  = \frac{1/Z_0 - 1/Z_1}{3}, \\[6pt]
  \mathbf{Y}_{abc} &=
  \begin{bmatrix}
    Y_s & Y_m & Y_m \\
    Y_m & Y_s & Y_m \\
    Y_m & Y_m & Y_s
  \end{bmatrix}.
\end{align}
The global admittance matrix is assembled from the standard bus-admittance
construction:
\begin{align}
  \mathbf{Y}^{abc}_{[3i{+}p,\;3i{+}q]} &\mathrel{+}= Y^{(b)}_{pq}
    + \tfrac{1}{2}\,jB_1^{(b)}\,\delta_{pq}, \\
  \mathbf{Y}^{abc}_{[3j{+}p,\;3j{+}q]} &\mathrel{+}= Y^{(b)}_{pq}
    + \tfrac{1}{2}\,jB_1^{(b)}\,\delta_{pq}, \\
  \mathbf{Y}^{abc}_{[3i{+}p,\;3j{+}q]} &\mathrel{-}= Y^{(b)}_{pq}, \\
  \mathbf{Y}^{abc}_{[3j{+}p,\;3i{+}q]} &\mathrel{-}= Y^{(b)}_{pq},
\end{align}
where $B_1^{(b)}$ is the positive-sequence line charging susceptance of branch
$b$, applied identically to each phase.

Per-phase bus shunt admittance is added on the diagonal when
\texttt{include\_shunts} is enabled:
$\mathbf{Y}^{abc}_{[3i{+}\phi,\;3i{+}\phi]} \mathrel{+}= g_\phi + jb_\phi$.

\paragraph{State vector and unknowns.}
The NR state vector is
\begin{align}
  \mathbf{x} = \bigl[
    \theta_{a,1},\theta_{b,1},\theta_{c,1},\ldots,
    V_{m,a,1},V_{m,b,1},V_{m,c,1},\ldots
  \bigr]^T.
\end{align}
\begin{itemize}
  \item \textbf{Slack bus:} All 3 phase angles and magnitudes are fixed
        (excluded from $\mathbf{x}$).
  \item \textbf{PQ bus:} Contributes $3\,\theta$ and $3\,V_m$ unknowns (6 per bus).
  \item \textbf{PV bus:} Contributes $3\,\theta$ unknowns (3 per bus); $V_m$ is fixed.
\end{itemize}
Denote the number of angle unknowns as $n_p$ and the number of magnitude
unknowns as $n_q$; the total system dimension is $n_p + n_q$.

\paragraph{Mismatch equations.}
At each non-slack phase node $r = 3i + \phi$:
\begin{align}
  P_r^{calc} &= \sum_{s=0}^{3n-1} |V_r||V_s|
    \bigl(G_{rs}\cos\theta_{rs} + B_{rs}\sin\theta_{rs}\bigr), \\
  Q_r^{calc} &= \sum_{s=0}^{3n-1} |V_r||V_s|
    \bigl(G_{rs}\sin\theta_{rs} - B_{rs}\cos\theta_{rs}\bigr),
\end{align}
with $\theta_{rs} = \theta_r - \theta_s$ and
$Y_{rs} = G_{rs} + jB_{rs}$.  The mismatch vector is
\begin{align}
  \Delta \mathbf{f} =
  \begin{bmatrix}
    P_r^{spec} - P_r^{calc} \\
    Q_r^{spec} - Q_r^{calc}
  \end{bmatrix},
  \qquad r \in \text{(angle or Vm unknowns)}.
\end{align}

\paragraph{Jacobian.}
The standard polar-form Jacobian entries for $r \neq s$ are:
\begin{align}
  \frac{\partial P_r}{\partial \theta_s} &=
    |V_r||V_s|\bigl(G_{rs}\sin\theta_{rs} - B_{rs}\cos\theta_{rs}\bigr), \\[3pt]
  \frac{\partial P_r}{\partial |V_s|} &=
    |V_r|\bigl(G_{rs}\cos\theta_{rs} + B_{rs}\sin\theta_{rs}\bigr), \\[3pt]
  \frac{\partial Q_r}{\partial \theta_s} &=
    -|V_r||V_s|\bigl(G_{rs}\cos\theta_{rs} + B_{rs}\sin\theta_{rs}\bigr), \\[3pt]
  \frac{\partial Q_r}{\partial |V_s|} &=
    |V_r|\bigl(G_{rs}\sin\theta_{rs} - B_{rs}\cos\theta_{rs}\bigr).
\end{align}
Diagonal (self) terms use the compact form:
\begin{align}
  \frac{\partial P_r}{\partial \theta_r} &= -Q_r^{calc} - |V_r|^2 B_{rr}, &
  \frac{\partial P_r}{\partial |V_r|} &= \frac{P_r^{calc}}{|V_r|} + |V_r| G_{rr}, \\[3pt]
  \frac{\partial Q_r}{\partial \theta_r} &= P_r^{calc} - |V_r|^2 G_{rr}, &
  \frac{\partial Q_r}{\partial |V_r|} &= \frac{Q_r^{calc}}{|V_r|} - |V_r| B_{rr}.
\end{align}

\paragraph{Update step.}
At each iteration the Newton correction is obtained by solving
\begin{align}
  \mathbf{J}\,\Delta\mathbf{x} = \Delta\mathbf{f},
\end{align}
using the shared sparse linear solver infrastructure
(\texttt{SparseLinearSolver}: KLU $>$ UMFPACK $>$ SparseLU fallback).
The update is
\begin{align}
  \theta_r^{(k+1)} &= \theta_r^{(k)} + \Delta\theta_r, \\
  |V_r|^{(k+1)} &= \max\!\bigl(|V_r|^{(k)} + \Delta|V_r|,\; 0.05\bigr),
\end{align}
with voltage clamping at $0.05\,\text{pu}$ to prevent divergence.

\paragraph{Convergence criterion.}
\begin{align}
  \varepsilon = \max_r \bigl|\Delta f_r\bigr| < \text{tol}.
\end{align}
Symbolic pattern analysis is performed once (iteration~0); subsequent iterations
reuse the sparsity pattern.

\paragraph{Post-convergence outputs.}
Per-phase branch current is recovered from the abc admittance:
\begin{align}
  I_b^{(\phi)} = \sum_{q=0}^{2} Y^{(b)}_{\phi q}
    \bigl(V_{from}^{(q)} - V_{to}^{(q)}\bigr),
\end{align}
producing sending-end power
$S_b^\phi = V_{from}^\phi \overline{I_b^\phi}$
and per-phase losses
$P_{loss}^\phi = |I_b^\phi|^2 \operatorname{Re}(Z_s)$,
$Q_{loss}^\phi = |I_b^\phi|^2 \operatorname{Im}(Z_s)$.

\subsubsection{Three-Phase Backward-Forward Sweep}
\label{sec:dpf-three-phase-bfs}

\paragraph{Note.} The BFS solver is retained as a legacy option for radial
single-source feeders.  The NR solver
(Section~\ref{sec:dpf-three-phase-nr}) is the recommended primary solver
because it supports meshed topologies, multiple generators, and hybrid AC/DC
networks.

\paragraph{Backward sweep.}
At each non-root bus $j$ in reverse BFS order, the per-phase injection current is
\begin{align}
  I_j^{\phi} = \overline{\frac{S_j^{net,\phi}}{V_j^{\phi}}},
  \qquad \phi \in \{a, b, c\}.
\end{align}
Branch currents are accumulated from children toward the root, preserving
three-phase coupling:
\begin{align}
  \mathbf{I}_{b(\text{parent})} \mathrel{+}= \mathbf{I}_{b(\text{child})}.
\end{align}

\paragraph{Forward sweep.}
For branch $b$ connecting parent bus $i$ to child bus $j$:
\begin{align}
  \mathbf{V}_j = \mathbf{V}_i - \mathbf{Z}_{abc}^{(b)}\, \mathbf{I}_b,
\end{align}
where $\mathbf{V}$, $\mathbf{I}$ are 3-vectors and
$\mathbf{Z}_{abc}^{(b)}$ is the $3\times3$ branch impedance.

\paragraph{Convergence.}
The stopping criterion extends to all phases:
\begin{align}
  \varepsilon = \max_{i,\,\phi} \bigl|V_i^{\phi,(k)} - V_i^{\phi,(k-1)}\bigr|.
\end{align}

\subsubsection{Per-Phase Branch Power and Losses}

After convergence, per-phase branch currents are recovered via the inverse of the
circulant impedance matrix:
\begin{align}
  \mathbf{Z}_{abc}^{-1} =
  \begin{bmatrix}
    Y_s & Y_m & Y_m \\
    Y_m & Y_s & Y_m \\
    Y_m & Y_m & Y_s
  \end{bmatrix},
  \qquad
  Y_s = \frac{1/Z_0 + 2/Z_1}{3}, \quad
  Y_m = \frac{1/Z_0 - 1/Z_1}{3},
\end{align}
and per-phase sending-end power:
\begin{align}
  S_b^{\phi} = V_{from(b)}^{\phi}\, \overline{I_b^{\phi}}.
\end{align}
Per-phase losses use the self-impedance component of $\mathbf{Z}_{abc}$:
\begin{align}
  P_{loss}^{\phi} = |I_b^{\phi}|^2 \operatorname{Re}(Z_s)\, S_{base},
  \qquad
  Q_{loss}^{\phi} = |I_b^{\phi}|^2 \operatorname{Im}(Z_s)\, S_{base}.
\end{align}
Total network losses are the sum across all branches and phases:
\begin{align}
  P_{loss}^{tot} = \sum_{\phi} P_{loss}^{\phi}, \qquad
  Q_{loss}^{tot} = \sum_{\phi} Q_{loss}^{\phi}.
\end{align}

\subsubsection{Voltage Unbalance Factor (VUF)}

The maximum voltage unbalance factor across all buses is reported as a quality
metric.  VUF is defined via the symmetrical-component decomposition of bus
voltages:
\begin{align}
  V_i^{+} &= \frac{1}{3}\bigl(V_i^a + a\,V_i^b + a^2\,V_i^c\bigr), \\
  V_i^{-} &= \frac{1}{3}\bigl(V_i^a + a^2\,V_i^b + a\,V_i^c\bigr), \\
  \mathrm{VUF}_i &= \frac{|V_i^{-}|}{|V_i^{+}|} \times 100\%.
\end{align}
A VUF of $0\%$ indicates a perfectly balanced system; the IEC 61000-2-2 standard
recommends $\mathrm{VUF} \le 2\%$ for low-voltage networks.

\subsubsection{Balanced-System Reduction}

When all three-phase loads are balanced ($S^a = S^b = S^c$) and all lines have
$Z_0 = Z_1$, both the NR and BFS abc-domain solvers produce per-phase results
identical to $1/3$ of the single-phase positive-sequence solution.  This
property is used for cross-validation of the three-phase implementations against
the existing single-phase solver.
